\section{Compositional Delegation Assurance}
\label{sec:compositional-delegation}

The trajectory semantics in Section~\ref{sec:trajectory-core} already includes multiple model, human, tool, and service actors; single-action execution is only its \(n=1\) case. Compositional delegation adds the propagation problem. The question is whether each assignment rested on a grant recognized under \(R\), each operation remained authorized, each cumulative and terminal condition held, and the record preserved the relevant basis across handoff, transformation, summarization, tool use, policy checking, and review.

This extension is better called \term{compositional delegation assurance} than simply multi-agent delegation assurance. Agents can receive delegated tasks, propose actions, and call tools. Tools may execute operations. Documents, contracts, policies, and memories may record grants, state constraints, or supply evidence; systems can also misclassify their contents. They don't acquire standing merely by entering a context. Treating all of them as agents would erase distinctions that delegation assurance needs to preserve. Treating them as nodes in a recorded propagation graph supplies one evaluation object, but not the normative hypergraph or execution trajectory.

In a compositional setting, the assurance representation is a graph with typed edges. Delegation edges record claims about grants and applicable constraints; retrieval, proposal, execution, and provenance edges record content flow, proposed operations, executed operations, or evidence. A joint grant or transition needs a guarded relation or hyperedge because several pairwise edges would falsely imply unilateral standing. The graph doesn't treat authority as a substance transmitted along edges, and a recorded edge doesn't by itself establish that the represented grant is valid.

A principal may authorize a coordinator agent to draft a procurement memo. The coordinator may ask a specialist agent to compare vendors, retrieve a contract, consult a policy, and prepare a tool-call proposal. The assurance question is whether each task assignment rested on a grant valid under \(R\), whether each proposed or executed operation fell within \(A_R(t)\), and whether the handoff record retained the governing sources, instruments, effective intervals, operation-specific standing, action class, scope, constraints, provenance, permission to subdelegate, information-flow permissions, persistence permissions, and non-delegable limits.

The normative and representational levels have to remain separate. Section~\ref{sec:delegation-beyond-responsiveness} defines \(A_R(t)\) from the operative normative state after valid transitions. Let \(s_t\) be the recorded authorization state, and let \(\widehat{A}(s_t)\) be the set of actions that state represents as permitted.

For an upstream grant \(g_{\mathrm{up}}\) held by a delegating source \(x\) and a proposed downstream grant \(g_{\mathrm{down}}\), grant-relative \term{non-amplification} requires
\[
  \operatorname{Scope}(g_{\mathrm{down}})
  \subseteq
  \operatorname{Conferrable}_{R}(x,g_{\mathrm{up}},t).
\]
The right-hand side is the region that \(R\) permits \(x\) to confer under that grant at \(t\); it needn't equal the operations \(x\) may personally execute. A separate valid authority-conferring event can supply a different basis for the downstream scope. Let \(\widehat{\operatorname{Scope}}(g)\) and \(\widehat{\operatorname{Conferrable}}(x,g,t)\) denote the corresponding regions represented in the record. The record-propagation condition requires every permission in
\[
  \widehat{\operatorname{Scope}}(g_{\mathrm{down}})
  \setminus
  \widehat{\operatorname{Conferrable}}(x,g_{\mathrm{up}},t)
\]
to be tied to a separately recorded grant or constraint change. This detects unexplained record expansion but doesn't establish that the recorded event is valid.

The record-soundness subclaim is \(\widehat{A}(s_t) \subseteq A_R(t)\), a relation that requires separate evidence under \(R\). An underinclusive record may preserve soundness while defeating a task-fulfilment, mandate-compliance, or record-completeness finding. Because authorization sets can be incomparable, none of these relations defines a single scalar amount of authority.

Adding an authority-conferring instrument can change \(\widehat{A}(s_t)\); it changes \(A_R(t)\) only if the instrument is valid and applicable under \(R\). Neither change rewrites an earlier grant or record state. A valid unrecorded grant can leave an operation authorized while creating a record gap at the corresponding assurance node and leaving a use-level conclusion unevaluable.

A summarization handoff has to retain the record of an applicable confidentiality constraint, and a downstream policy-enforcement point has to apply that constraint to the proposed operation. A coordinator has to keep \mention{draft an internal comparison} distinct from \mention{contact the vendor}. A tool wrapper has to require the approval specified by the governing rule rather than treat a natural-language rationale as approval. A retrieved contract clause may supply evidence of an applicable constraint, subject to validity and adjudication, but its presence in context doesn't make it an operation to execute.

The three structures now do distinct work. \(H_R(t)\) supplies the normative relations against which grants and events are assessed. \(\widehat G_B\) shows which asserted grants, constraints, and classifications were propagated. \(\pi_B\) records the ordered events and state changes within the boundary. OpenAI's report of fragmented-token reconstruction illustrates why the last object matters, but its narrative and monitoring proposal don't establish the trajectory condition or validate a monitor \parencite{openai2026longHorizonSafety}.

Compositional delegation assurance adds several requirements to the single-action pattern (Table~\ref{tab:delegation-evidence}). First, the evidence has to support that the governing grant permits subdelegation. Second, the normative authorization path has to satisfy non-amplification. Third, every handoff record has to satisfy the record-propagation condition, with every added permission tied to a separately recorded authority-conferring event and with the event's validity and applicability under \(R\) assessed separately for record soundness.

Fourth, the record has to retain every still-applicable constraint through summarization, routing, and memory updates. Fifth, provenance has to remain linked to instructions, documents, tool results, and policy text. Sixth, unresolved conflicts have to be escalated or recorded as conflicts rather than converted into an unqualified action proposal. Seventh, a trajectory claim has to declare \(B\), retain ordered operations and intermediate states, and evaluate applicable local, prefix, cumulative, and terminal constraints separately. Finally, the audit record has to reconstruct both the normative basis asserted for the action and the sequence of classifications, proposals, gate decisions, monitor interventions, and executions from the principal's grant to the contested action or bounded trajectory.

Several failure levels have to remain separate. An intermediate record may omit a grant or still-applicable constraint during handoff. The omission is a recorded-authorization propagation failure and defeats the corresponding component or audit subclaim. It creates a record gap at that node and may defeat the subclaim or leave the declared use-level conclusion unevaluable, depending on the claim tree, without making the normative proposition that the executed operation fell within \(A_R(t)\) false. An unsupported permission introduced into \(\widehat{A}(s_t)\) is a record-soundness failure even if a later gate blocks the proposed operation; an executed operation outside \(A_R(t)\) defeats the action-authorization claim. A trajectory can also satisfy every applicable step-local condition while violating an independently specified cumulative or terminal constraint. That defeats the trajectory claim without retroactively making each component operation locally unauthorized.

Separately, accumulated context may destabilize behaviour even when the record remains intact. A changed output motivates a contextual-reliability hypothesis. A trace may independently support an item-level classification--adjudication misfit concerning status, standing, priority, or scope, but no single item establishes a contextual-reliability failure over an evaluation window. The projective claim requires target-matched repeated evidence. Path tests should include matched non-controlling loads and repeated samples rather than infer a dropped or misclassified grant or constraint from changed output alone \parencite{zhang2026illusionRobustness}.

This separation also governs automated-evaluator validation. A final-transcript score may miss an unsupported permission introduced at an intermediate node or a cumulative constraint violated only by the sequence. An automated evaluator for compositional delegation may be assigned to assess whether a subagent was permitted to receive the relevant information, whether it had a valid grant for the assigned action, whether the proposed operation belonged to \(\widehat{A}(s_t)\), whether the executed operation belonged to \(A_R(t)\), whether local, prefix, cumulative, and terminal constraints were satisfied over \(\pi_B\), whether constraints remained recorded across summaries, and whether the trace supports third-party reconstruction.

That label is evidence about the assigned claims, not a grant, authorization decision, or liability finding. If a decision rule gives it operational effect, the assurance graph should record who authorized its use, the evaluator's scope, and the rule.

\begin{figure}[tbp]
\centering
\begin{tikzpicture}[
  font=\scriptsize,
  gnode/.style={draw=black!55, thick, rounded corners=2pt, align=center, inner sep=2pt, minimum height=0.58cm, text width=1.4cm, fill=black!3},
  flow/.style={-{Latex[length=1.6mm]}, thick, draw=black!60},
  norm/.style={-{Latex[length=1.6mm]}, thick, dashed, draw=linkmaroon},
  panel/.style={draw=black!35, rounded corners=2pt}
]
% Panel (a): normative hypergraph
\draw[panel] (-0.8,1.4) rectangle (3.3,-3.8);
\node[font=\footnotesize, align=center] at (1.25,0.98) {(a) normative hypergraph\\\(H_R(t)\)};
\node[gnode] (p) at (0,0.1) {Procurement\\principal};
\node[gnode] (co) at (2.5,0.1) {Compliance\\officer};
\node[gnode, fill=linkmaroon!7] (guard) at (1.25,-1.35) {Guarded\\joint grant};
\node[gnode] (coord) at (1.25,-2.75) {Coordinator\\mandate};
\draw[norm] (p) -- (guard);
\draw[norm] (co) -- (guard);
\draw[norm] (guard) -- node[right, font=\tiny]{valid if both approve} (coord);

% Panel (b): recorded graph
\begin{scope}[xshift=4.35cm]
\draw[panel] (-0.55,1.4) rectangle (3.75,-3.8);
\node[font=\footnotesize, align=center] at (1.6,0.98) {(b) recorded graph\\\(\widehat G_B\)};
\node[gnode] (rp) at (0.15,-0.15) {Principal\\record};
\node[gnode] (rc) at (1.65,-1.35) {Coordinator\\record};
\node[gnode] (rs) at (3.05,-2.35) {Specialist\\record};
\node[gnode] (rg) at (1.65,-3.05) {Gate\\record};
\draw[flow] (rp) -- node[above right, font=\tiny]{handoff} (rc);
\draw[flow] (rc) -- node[above right, font=\tiny]{summary} (rs);
\draw[flow] (rs) -- node[below, font=\tiny]{proposal} (rg);
\end{scope}

% Panel (c): execution trajectory
\begin{scope}[xshift=9.15cm]
\draw[panel] (-0.65,1.4) rectangle (3.65,-3.8);
\node[font=\footnotesize, align=center] at (1.5,0.98) {(c) execution trajectory\\\(\pi_B\)};
\node[gnode] (e1) at (0,-0.25) {\(e_1\) model\\reads};
\node[gnode] (e2) at (1.5,-1.55) {\(e_2\) human\\approves};
\node[gnode] (e3) at (3.0,-2.65) {\(e_3\) tool\\sends};
\draw[flow] (e1) -- node[above right, font=\tiny]{\(z_0\!\to z_1\)} (e2);
\draw[flow] (e2) -- node[above right, font=\tiny]{\(z_1\!\to z_2\)} (e3);
\node[font=\tiny, align=center, text=linkmaroon] at (1.45,-3.25) {local, prefix, cumulative, and terminal\\conditions evaluated over history};
\end{scope}
\end{tikzpicture}
\caption{Three non-equivalent compositional structures. Dashed maroon edges in (a) are normative relations under \(R\), solid arrows in (b) are recorded propagation, and solid arrows in (c) order executed events and state changes. A pairwise recorded graph can't replace the guarded hyperedge, and neither graph establishes what occurred.}
\label{fig:delegation-triptych}
\end{figure}
\FloatBarrier

The three programmes in Section~\ref{sec:comparative-test} extend to the structures in Figure~\ref{fig:delegation-triptych} without merging them. Randomized local discrimination can assign packages that represent the same handoffs and proposed call while varying whether an intermediate record retains an applicable constraint. A second local contrast holds the proposed call and technical capabilities constant while varying whether the package represents a valid subdelegation.

Predictive ablation compares held-out prediction from a typed \(\widehat G_B\) and event representation against weaker feature sets, using an allowlisted view from which reference and oracle fields are absent. Reviewer reconstruction holds \(H_R(t)\) and \(\pi_B\) fixed while independently randomizing intact or degraded versions of \(\widehat G_B\); neither view contains the hidden reference. The unit in the latter programme is a reviewer judgment on a trace, not the model or the trajectory.

A factorial family varies contextual status while holding source standing fixed and varies source standing while holding contextual status fixed. Another holds the ordered events and local permissions fixed while varying whether \(R^{\mathrm{stip}}\) contains a cumulative or terminal condition that the trajectory violates. A matched placebo places the same restriction in quoted or otherwise non-operative text. Static transcript cases test local discrimination or reviewer reconstruction; a live tool-mediated harness is a separate, horizon-indexed projective target.

Purposive fixtures can expose level errors and local failure modes. They don't estimate real-world prevalence. These tests can show whether recorded classifications, observable decisions, predictions, or reviewer judgments respond to the tested contrasts. They don't by themselves establish internal representation, legal validity, or performance on other paths, systems, regimes, reviewer populations, deployments, or times.
